% ==============================================================
% Projektive Strukturtheorie (PST)
% Dokument PST-6-R: Physikalische Rekonstruktion
% Version 0.1
% Kompiliert mit: pdfLaTeX
% ==============================================================

\documentclass[11pt,a4paper]{article}

% ---------- Grundpakete ----------
\usepackage[utf8]{inputenc}
\usepackage[T1]{fontenc}
\usepackage[german]{babel}
\usepackage{amsmath,amssymb,amsthm}
\usepackage{geometry}
\usepackage{parskip}
\usepackage{enumitem}
\usepackage{booktabs}
\usepackage{xcolor}
\usepackage{hyperref}

% ---------- Layout ----------
\geometry{margin=2.5cm}
\setlist{itemsep=0.2em, topsep=0.2em}
\hypersetup{
    colorlinks=true,
    linkcolor=blue,
    urlcolor=blue,
    pdftitle={PST-6-R: Physikalische Rekonstruktion -- Version 0.1},
    pdfauthor={Forschungsprogramm Ontoph\"{a}nomenalismus}
}

% ---------- Theorem-Umgebungen ----------
\theoremstyle{plain}
\newtheorem{theorem}{Theorem}[section]
\newtheorem{proposition}{Proposition}[section]
\newtheorem{lemma}{Lemma}[section]

\theoremstyle{definition}
\newtheorem{definition}{Definition}[section]
\newtheorem{remark}{Bemerkung}[section]
\newtheorem{methodennote}{Methodennote}[section]
\newtheorem{hypothese}{Arbeitshypothese}[section]
\newtheorem{programm}{Forschungsprogramm}[section]

% Neue Kommandos
\newcommand{\OPP}{\mathcal{M}_{\mathrm{OPP}}}
\newcommand{\deff}{d_{\mathrm{eff}}}
\newcommand{\Proj}{\mathcal{P}}
\newcommand{\Lp}{\ell_P}

% ---------- Titel ----------
\title{
  \textbf{Projektive Strukturtheorie (PST)}\\[0.4em]
  \textbf{Dokument PST-6-R: Physikalische Rekonstruktion}\\[0.4em]
  {\large Lorentz-Signatur, Eichgruppen, Quantisierung,
  Naturkonstanten und das Fernziel}\\[0.5em]
  {\normalsize Version~0.1}
}
\author{
  \textbf{Forschungsprogramm Ontoph\"{a}nomenalismus}\\
  \small Siehe TRANSPARENCY f\"{u}r methodologische Urheberschaft
  und KI-Einsatz.
}
\date{Juli 2026}

\begin{document}

\maketitle

\begin{abstract}
Dieses Dokument entwickelt die physikalische Rekonstruktion der PST ---
den Schritt von der emergenten Raumzeit zur vollst\"{a}ndigen Physik.

PST-6-R ist das ambiti\"{o}seste und am wenigsten ausgearbeitete
der PST-Dokumente. Es kartografiert das Terrain das die PST noch
zu erschlie\ss{}en hat: Lorentz-Signatur, Eichgruppen, Quantisierung
und Naturkonstanten.

Gem\"{a}\ss{} MKO sind alle Aussagen in diesem Dokument explizit als
Arbeitshypothesen (Ebene~5) gekennzeichnet. PST-6-R ist kein Ergebnisbericht
--- es ist ein Forschungsprogramm.

Es gelten die methodologischen Konventionen gem\"{a}\ss{} MKO.
\end{abstract}

\tableofcontents
\newpage

%==============================================================
\section{Einordnung: Das letzte Glied der Kette}
%==============================================================

\begin{methodennote}[Status von PST-6-R --- MKO Ebene~5]
PST-6-R ist das spekulativste der PST-Dokumente. W\"{a}hrend PST-1-F
bis PST-4-G gesicherte Grundlagen und numerische Evidenz beschreiben,
kartografiert PST-6-R das noch unerschlossene Terrain.

Alle Hypothesen in diesem Dokument sind MKO Ebene~5: empirisch testbar
in Prinzip, aber noch nicht getestet. Dieses Dokument ist eine
Einladung zur weiteren Forschung, kein Abschluss.
\end{methodennote}

Die Rekonstruktionskette der PST:
\[
  \underbrace{\text{Existenz} \rightarrow \OPP}_{\text{PST-1-F}}
  \rightarrow
  \underbrace{\mathbf{K} \rightarrow \deff \approx 4}_{\text{PST-2-M, PST-4-G}}
  \rightarrow
  \underbrace{\text{Raumzeit}}_{\text{PST-3-P}}
  \rightarrow
  \underbrace{\text{Lorentz} \rightarrow \text{Eich} \rightarrow \hbar
  \rightarrow \text{Physik}}_{\text{PST-6-R (offen)}}
\]

%==============================================================
\section{Lorentz-Signatur: Das n\"{a}chste gro\ss{}e Ziel}
%==============================================================

\begin{methodennote}[Wo wir stehen]
Die PST hat numerisch $\deff \approx 4$ etabliert. Die emergierte
Symmetriegruppe ist bisher $SO(4)$ --- euklidische 4D-Symmetrie.
Die physikalische Raumzeit hat jedoch $SO(1,3)$ --- Lorentz-Symmetrie
mit Signatur $(+,-,-,-)$.

Der Unterschied ist fundamental: $SO(4)$ hat kein ausgezeichnetes
Zeitintervall; $SO(1,3)$ hat Zeitpfeile, Kausalit\"{a}t und Lichtkegelstruktur.
\end{methodennote}

\begin{hypothese}[Wick-Rotation als projektive Eigenschaft --- MKO Ebene~5]
Die \"{U}bertragung von euklidischer zu Lorentz-Signatur k\"{o}nnte durch
eine \emph{Wick-Rotation} des Projektionsoperators beschrieben werden:
\[
  t \;\rightarrow\; -i\tau
\]
wobei $t$ die projizierte Zeitkoordinate und $\tau$ die euklidische Zeit
bezeichnet. In der QFT ist die Wick-Rotation ein Standardtrick;
in der PST k\"{o}nnte sie eine tiefere Bedeutung haben: Die Lorentz-Signatur
emergiert durch Komplexifizierung des Projektionsoperators.
\end{hypothese}

\begin{hypothese}[Imagin\"{a}re Eigenwerte --- MKO Ebene~5]
Eine alternative Hypothese: Falls die Gram-Matrix $\mathbf{K}$ komplexwertige
Eintr\"{a}ge erlaubt (etwa durch analytische Fortsetzung der
Relationsst\"{a}rke), k\"{o}nnten imagin\"{a}re Eigenwerte die negativen
Vorzeichen der Lorentz-Signatur erzeugen:
\[
  \lambda_{\text{Zeit}} = i\mu, \quad \mu > 0
  \;\Rightarrow\;
  \sqrt{\lambda_{\text{Zeit}}} = i\sqrt{\mu}
  \;\Rightarrow\;
  g_{00} = -\mu < 0.
\]
Dies ist ein konkreter mathematischer Ansatz, der weiter gepr\"{u}ft werden muss.
\end{hypothese}

\begin{programm}[Lorentz-Programm]
\textbf{Ziel:} Mathematische Herleitung der Signatur $(+,-,-,-)$ aus
dem Projektionsprozess.

\textbf{Strategie:}
\begin{enumerate}
  \item Komplexifizierung der Gram-Matrix und Analyse der Eigenwertstruktur.
  \item Vergleich Wick-Rotation vs.\ imagin\"{a}re Eigenwerte.
  \item Nachweis dass $SO(1,3)$ die einzige stabile Signatur f\"{u}r
        $\deff = 4$ unter physikalischen Randbedingungen ist.
\end{enumerate}

\textbf{Kriterium f\"{u}r Erfolg:} $SO(1,3)$ emergiert rangneutral
--- nicht weil es gesucht wurde, sondern weil es die einzige
stabile Signatur ist.
\end{programm}

%==============================================================
\section{Eichgruppen: Die Materiestruktur}
%==============================================================

\begin{hypothese}[Eichgruppen als Projektiv-Symmetrien --- MKO Ebene~5]
Die Eichgruppen des Standardmodells
$G_{\text{SM}} = SU(3) \times SU(2) \times U(1)$
sind Untergruppen der Symmetriegruppe des OPP, die nach der physikalischen
Projektion erhalten bleiben:
\[
  \Proj^\ast(G_{\OPP}) \;\supseteq\; G_{\text{SM}} \times SO(1,3).
\]
Materie und Wechselwirkungen sind damit keine separaten ontologischen
Einheiten --- sie sind Aspekte der Projektionssymmetrie.
\end{hypothese}

\begin{remark}[Warum gerade $SU(3) \times SU(2) \times U(1)$?]
Die PST stellt die tiefste Frage zur Materiestruktur: Warum haben wir
genau diese Eichgruppen und nicht andere?

Kandidatenantworten:
\begin{enumerate}
  \item \textbf{Stabilit\"{a}t:} $SU(3) \times SU(2) \times U(1)$ ist
        die einzige Eichgruppe die unter der physikalischen Projektion
        stabil ist.
  \item \textbf{Minimalit\"{a}t:} Es ist die kleinste Gruppe die alle
        bekannten Teilchen klassifiziert.
  \item \textbf{Kompatibilit\"{a}t:} Sie ist die einzige Gruppe die mit
        $SO(1,3)$ und $\deff = 4$ kompatibel ist.
\end{enumerate}
Welche Antwort zutrifft, ist eine offene Forschungsfrage.
\end{remark}

\begin{programm}[Eichgruppen-Programm]
\textbf{Ziel:} Ableitung von $SU(3) \times SU(2) \times U(1)$ aus
Eigenschaften des Projektionsoperators $\Proj^\ast$.

\textbf{Strategie:}
\begin{enumerate}
  \item Charakterisierung der Symmetriegruppen verschiedener
        Projektionskandidaten.
  \item Pr\"{u}fung ob Diffusion Maps mit $t \approx 9$--$10$ spezifische
        Symmetrien bevorzugt.
  \item Verbindung zur Connes'schen nicht-kommutativen Geometrie
        (spektrale Tripel kodieren Eichstrukturen).
\end{enumerate}
\end{programm}

%==============================================================
\section{Quantisierung: Diskretheit aus Projektion}
%==============================================================

\begin{hypothese}[Quantisierung als Projektionsartefakt --- MKO Ebene~5]
Quantisierung entsteht durch die endliche Aufl\"{o}sung des
Projektionsoperators. Die Informationsschwelle $I_0$ (PST-2-M)
entspricht dem Planck'schen Wirkungsquantum:
\[
  I_0 \;\leftrightarrow\; \hbar.
\]
Unterhalb von $I_0$ gibt es keine projektive Bestimmtheit ---
keine unterscheidbaren Quantenzust\"{a}nde.
\end{hypothese}

\begin{remark}[Drei Ebenen der Quantisierung]
Die PST-Hypothese ber\"{u}hrt drei verschiedene Aspekte der Quantisierung:
\begin{enumerate}
  \item \textbf{Diskretheit von Energie:} $E = n\hbar\omega$ ---
        k\"{o}nnte aus der diskreten Spektralstruktur der Gram-Matrix folgen.
  \item \textbf{Unsch\"{a}rfe:} $\Delta x \cdot \Delta p \geq \hbar/2$
        --- k\"{o}nnte aus der Nicht-Injektivit\"{a}t des Projektionsoperators
        folgen (PST-3-P).
  \item \textbf{Spin und Statistik:} Das Spin-Statistik-Theorem k\"{o}nnte
        aus Eigenschaften der projizierten Symmetriegruppe $SO(1,3)$ folgen.
\end{enumerate}
\end{remark}

\begin{hypothese}[Schr\"{o}dinger-Gleichung aus Projektion --- MKO Ebene~5]
Die Schr\"{o}dinger-Gleichung
\[
  i\hbar\, \frac{\partial\psi}{\partial t} = \hat{H}\psi
\]
k\"{o}nnte eine projektive Eigenschaft sein: die Beschreibung wie sich
ein projizierter Zustand $\psi$ im projizierten Zeitparameter $t$
ver\"{a}ndert. Die Hamilton-Operator $\hat{H}$ w\"{a}re dann eine
Eigenschaft des Projektionsoperators, nicht des OPP selbst.

Vorsicht: Diese Hypothese ist sehr spekulativ. Sie w\"{u}rde bedeuten dass
die gesamte Quantenmechanik aus der Projektion folgt --- ein
au\ss{}ergew\"{o}hnlicher Anspruch der au\ss{}ergew\"{o}hnlicher Evidenz
bed\"{u}rfte.
\end{hypothese}

%==============================================================
\section{Naturkonstanten: Die h\"{a}rteste Pr\"{u}fung}
%==============================================================

\begin{methodennote}[Naturkonstanten als Pr\"{u}fstein]
Naturkonstanten sind die pr\"{a}zisesten Messgr\"{o}\ss{}en der Physik.
Die Feinstrukturkonstante $\alpha \approx 1/137.036$ ist auf viele
Stellen bekannt. Eine Theorie die $\alpha$ aus Grundprinzipien ableitet,
w\"{a}re ein au\ss{}ergew\"{o}hnlicher Erfolg.

Die PST macht bisher keine Vorhersage \"{u}ber Naturkonstanten.
Dieser Abschnitt kartografiert das Terrain.
\end{methodennote}

\begin{hypothese}[Planck-L\"{a}nge aus Existenzgrenze --- MKO Ebene~5]
Die Planck-L\"{a}nge
\[
  \Lp = \sqrt{\frac{\hbar G}{c^3}} \approx 1.616 \times 10^{-35}\,\text{m}
\]
k\"{o}nnte mit der Informationsschwelle $I_0$ zusammenh\"{a}ngen
(PST-2-M, Sekis Beitrag): Unterhalb von $\Lp$ gibt es keine projektive
Bestimmtheit. Das w\"{u}rde $\Lp$ aus OPP-Eigenschaften ableiten.

Konkret: Falls $I_0$ aus der Spektralstruktur des OPP hergeleitet werden
kann, und falls $I_0 \leftrightarrow \hbar/\ell_P^2$ gilt, folgt $\ell_P$
aus den OPP-Eigenschaften.
\end{hypothese}

\begin{hypothese}[Lichtgeschwindigkeit als Projektionsparameter ---
MKO Ebene~5]
Die Lichtgeschwindigkeit $c$ k\"{o}nnte der Konversionsparameter
zwischen der zeitartigen und den raumartigen Dimensionen der
projizierten Raumzeit sein --- eine Eigenschaft des Projektionsoperators
$\Proj^\ast$, nicht des OPP selbst. In diesem Sinne w\"{a}re $c$ kein
fundamentaler OPP-Parameter, sondern eine Eigenschaft der Projektion.
\end{hypothese}

\begin{remark}[Feinstrukturkonstante: noch offen]
Die Feinstrukturkonstante $\alpha \approx 1/137$ ist in der PST bisher
vollst\"{a}ndig offen. Richard Feynman beschrieb $\alpha$ als eines der
gr\"{o}\ss{}ten R\"{a}tsel der Physik. Die PST stellt die Frage ob
$\alpha$ aus Eigenschaften des Projektionsoperators folgt --- aber
hat noch keine Antwort.
\end{remark}

%==============================================================
\section{Dynamik: Wie bewegt sich etwas im statischen OPP?}
%==============================================================

\begin{methodennote}[Das Dynamikproblem]
Der OPP ist statisch --- er enth\"{a}lt alle Zust\"{a}nde zeitlos und
simultan (vgl.\ OPP). Wie entsteht daraus die Dynamik des physikalischen
Universums?
\end{methodennote}

\begin{hypothese}[Dynamik als Beobachterpfad --- MKO Ebene~5]
Physikalische Dynamik ist kein Prozess im OPP selbst. Sie ist die
\emph{Sequenz von Projektionen} die ein Beobachter-Attraktor durchl\"{a}uft:
\[
  \gamma\colon \mathbb{R} \rightarrow \OPP,
  \quad
  t \mapsto \gamma(t)
\]
wobei $t$ der \emph{projizierte} Zeitparameter ist --- kein OPP-Parameter.
Dynamik ist damit eine Eigenschaft der Beobachter-Projektion, nicht
des OPP. Dies ist konsistent mit der Blockuniversum-Analogie in OPP
und dem Pausen-Argument in DPR.
\end{hypothese}

\begin{hypothese}[Kausalit\"{a}t aus Projektionsreihenfolge ---
MKO Ebene~5]
Kausalit\"{a}t --- die Tatsache dass Ursachen vor Wirkungen kommen ---
k\"{o}nnte aus der Reihenfolge der Projektionen folgen: Ein Zustand
$\gamma(t_1)$ verursacht $\gamma(t_2)$ genau dann, wenn die Projektion
von $\gamma(t_1)$ in $\gamma(t_2)$ hineinreicht (informationale Ursache),
nicht umgekehrt. Die Zeitrichtung w\"{a}re dann eine Eigenschaft der
Projektionssequenz.
\end{hypothese}

%==============================================================
\section{Verbindung zu anderen Fundamentaltheorien}
%==============================================================

\begin{remark}[Allgemeine Relativit\"{a}tstheorie (ART)]
Die ART beschreibt Gravitation als Kr\"{u}mmung der Raumzeit. In der PST
w\"{a}re Raumzeitkr\"{u}mmung eine Eigenschaft des Projektionsoperators:
Kr\"{u}mmung entsteht wenn der Projektionsoperator in verschiedenen
Regionen des OPP verschiedene effektive Dimensionen liefert.

Die Einstein-Feldgleichungen
\[
  G_{\mu\nu} + \Lambda g_{\mu\nu} = \frac{8\pi G}{c^4} T_{\mu\nu}
\]
k\"{o}nnten aus Eigenschaften des Projektionsspektrums emergieren. Dies
ist eine sehr spekulative Hypothese --- aber konsistent mit der
PST-Grundidee.
\end{remark}

\begin{remark}[Kosmologische Konstante $\Lambda$]
Die kosmologische Konstante $\Lambda$ ist ein offenes Problem der
Physik: Quantenfeldtheorie sagt einen Wert vorher der um $10^{120}$
gr\"{o}\ss{}er ist als beobachtet (Feinabstimmungsproblem). In der PST
k\"{o}nnte $\Lambda$ eine Eigenschaft des OPP selbst sein --- der
\glqq{}Grundzustand\grqq{} der Projektion. Falls der OPP eine bestimmte
Kohärenz-Dichte hat, k\"{o}nnte $\Lambda$ daraus folgen. Diese Hypothese
ist sehr spekulativ.
\end{remark}

%==============================================================
\section{Das Fernziel: Existenz $\Rightarrow$ Universum}
%==============================================================

\begin{remark}[Was erreicht sein m\"{u}sste]
Das Fernziel der PST ist erreicht wenn alle folgenden Schritte
mathematisch notwendig sind:
\[
  \text{Existenz}
  \xrightarrow{(1)} \OPP
  \xrightarrow{(2)} \Proj^\ast
  \xrightarrow{(3)} \text{Raumzeit } 3{+}1
  \xrightarrow{(4)} SO(1,3)
  \xrightarrow{(5)} G_{\text{SM}}
  \xrightarrow{(6)} \hbar, c, G
  \xrightarrow{(7)} \text{Teilchen, Felder}
  \xrightarrow{(8)} \text{Universum}
\]
Schritt (1) und Teile von (2) sind etabliert. Schritt (3) ist numerisch
best\"{a}tigt. Schritte (4)--(8) sind offen.
\end{remark}

\begin{remark}[Warum das Fernziel trotzdem verfolgt werden sollte]
Das Fernziel ist m\"{o}glicherweise nicht vollst\"{a}ndig erreichbar ---
G\"{o}delsche Unvollst\"{a}ndigkeit und physikalische Komplexit\"{a}t
setzen Grenzen. Aber jeder erreichte Schritt ist ein echter
wissenschaftlicher Fortschritt. Und das Fernziel gibt dem Programm
Richtung und Koh\"{a}renz.

Vergleich: Die vollst\"{a}ndige Vereinheitlichung der Physik
(Theory of Everything) ist m\"{o}glicherweise ebenfalls nicht erreichbar.
Trotzdem hat die Suche danach Quantenmechanik, Relativit\"{a}tstheorie
und das Standardmodell hervorgebracht.
\end{remark}

%==============================================================
\section{Offene Fragen der physikalischen Rekonstruktion}
%==============================================================

\begin{enumerate}
  \item Wie emergiert die Lorentz-Signatur $(+,-,-,-)$ aus einer
        reellen oder komplex-erweiterten Gram-Matrix?
  \item Welche Symmetrien des OPP bleiben unter $\Proj^\ast$ erhalten,
        und warum gerade $G_{\text{SM}} \times SO(1,3)$?
  \item Ist $I_0 \leftrightarrow \hbar$ mathematisch pr\"{a}zisierbar,
        und folgt daraus die Planck-L\"{a}nge?
  \item Wie beschreibt man Dynamik in einem statischen OPP ---
        und folgt die Schr\"{o}dinger-Gleichung aus der Projektion?
  \item K\"{o}nnen die Einstein-Feldgleichungen aus Eigenschaften des
        Projektionsspektrums emergieren?
  \item Gibt es eine Verbindung zwischen der kosmologischen Konstante
        $\Lambda$ und der Koh\"{a}renz-Dichte des OPP?
  \item Wie h\"{a}ngt die Feinstrukturkonstante $\alpha$ mit
        OPP-Eigenschaften oder $\Proj^\ast$ zusammen?
  \item Ist die PST mit Loop-Quanten-Gravitation oder
        nicht-kommutativer Geometrie (Connes) formal \"{a}quivalent?
\end{enumerate}

%==============================================================
\section{Schnittstellen}
%==============================================================

\begin{itemize}
  \item \textbf{PST-3-P:} Die Projektionstheorie aus PST-3-P ist die
        Grundlage f\"{u}r alle physikalischen Rekonstruktionen in PST-6-R.
        Insbesondere Symmetriebrechung, Beobachter als Operatoren und
        $\Proj^\ast$.
  \item \textbf{PST-5-S:} Die priorisierte Forschungsagenda aus PST-5-S
        bestimmt die Reihenfolge der Aufgaben f\"{u}r PST-6-R.
  \item \textbf{MGO:} Kategorientheorie, Operatoralgebren und
        nicht-kommutative Geometrie (MGO) sind die mathematischen
        Werkzeuge f\"{u}r PST-6-R.
  \item \textbf{DPR:} Das Pausen-Argument und die RdE aus DPR sind
        konzeptuelle Grundlagen f\"{u}r die Dynamik-Hypothese.
  \item \textbf{OPP:} Die Blockuniversum-Analogie aus OPP ist
        konsistent mit der Dynamik-als-Beobachterpfad-Hypothese.
  \item \textbf{DFD:} Die Planck-Schranke als Existenzgrenze
        (PST-2-M, Sekis Beitrag) verbindet DFD mit PST-6-R:
        $I_0 \leftrightarrow \hbar \leftrightarrow \mathcal{D}(E)_{\min}$.
\end{itemize}

\end{document}

